% !TEX root = main.tex

\section{Reachability effect}\label{sect:typesystem}


A {\em reachability graph $\graph$} is a directed graph with a finite set of nodes 
labelled by distinct variables. If in $\graph$ there is a path from
a node labelled $\x$ to one labelled $\y$ then the location referred by $\y$ could also be reached from $\x$.
We represent $\graph$ with the set of its edges, $(\x,\y)$, where $\x$ is the label of the source node and
$\y$ the one of the destination of the edge. $\Vars(\graph)$ is the set of variables labelling the nodes of $\graph$.

The reachability judgments have the shapes\[
\IsWFValGr\Gamma\ve\tau\VarSet\qquad\IsWFExpGr\Gamma\e\tau\graph
\]
where $\Gamma$ is a \emph{type {environment}}, written $\VarType{\x_1}{\tau_1},\ldots,\VarType{\x_n}{\tau_n}$, $\VarSet$ a set of variables and
$\graph$ is a reachability graph.  
The labels of the nodes of $\graph$ are variables in $\dom{\Gamma}$, plus a distinguished variable $\res$  for ``result''. 
In the judgment for values, $\VarSet$ is the set of variables reachable from $\ve$.
In the judgment for expressions, the 
meaning is that $\graph$ represents the ``reachability''??  among the free
variables of $\e$ possibly introduced by the evaluation of the expression, and
the variables  $\x$ such that there is a path from $\x$ to $\res$ are the ones that will be
possibly refer to the result of the expression.  

For reachability graph we use the following notations:
\begin{itemize}
\item $\emptygraph$ for the empty graph
  \item $\AddEdge\x\y\graph=\graph\cup\{(\x,\y)\}$ adds the edge from the node labelled $\x$ to the one labelled $\y$. 
  \item $\SumGr{\graph_1}{\graph_2}=\graph_1\cup\graph_2$ 
  \item $\RemoveGr\graph\VarSet=\{(\y,\z)\in\graph\mid\y\not\in\VarSet\ \wedge\ \z\not\in\VarSet)\}\cup\{(\y,\z)\in\graph\mid\x\in\VarSet\ \wedge\ (\y,\x)\in\graph\ \wedge\ (\x,\z)\in\graph\}$ removes the node with labels in $\VarSet$ from $\graph$ propagating their reachability relation
  \item $\SubstGr\graph\VarSet\x=\{(\y,\z)\in\graph\mid(\x,\z)\in\graph\ \wedge\ \y\in\VarSet\ \wedge\ \y\neq\z\}\cup\{(\z,\y)\in\graph\mid(\z,\x)\in\graph\ \wedge\ \y\in\VarSet\ \wedge\ \y\neq\z\}$ we assume  $\VarSet\cap\Vars(G)=\emptyset$.
  \end{itemize}
  The type of functions should be decorated by the effect of the execution of the function body. Say $\graph$ is the effect of the typing of the function body. 
  \begin{itemize}
  \item The effects on the free variables are given by $\graph$.
  \item For the effects on the ``actual'' parameter we decorate $\graph$ with the name of the ``formal'' parameter (questa \`e la dipendenza)
 \end{itemize}
Change to the definition of types:

\begin{grammatica} 
\produzione{\sigma,\tau}{\RefTV{\VarSet}{\tau}\mid\funTypeGr{\tau_1}{\tau_2}{\x}{\graph}\mid\dots}{}
\end{grammatica}
\\
Predicate $\refType{\tau}$ holds if 
\begin{itemize}
\item $\tau=\RefTV{\VarSet}{\tau'}$ for some $\tau'$ and $\VarSet$ or
\item $\tau={\PairT{\tau_1}{\tau_2}}$ and either $\refType{\tau_1}$ or $\refType{\tau_2}$.
 \end{itemize}
 
\begin{figure}
\begin{math}
\begin{array}{l}
\NamedRule{t-var}{}
{ \IsWFValGr{\VarType{\x}{\tau}}{\x}{\tau}{\VarSet} } 
{\VarSet=\begin{cases}
\{\x\} &  \refType{\tau}\\
\emptyset  & \text{otherwise}
\end{cases}}
\BigSpace
\NamedRule{te-basic}{}
{
\IsWFValGr{\emptyset}{\bv}{{\BT}}{\emptyset}
}
{}
\\[3ex]
\NamedRule{te-pair}{\IsWFValGr{\Gamma_1}{\ve_1}{{\tau_1}}{\VarSet_1}\BigSpace\IsWFValGr{\Gamma_2}{\ve_2}{{\tau_2}}{\VarSet_2}}
{
\IsWFValGr{\Gamma_1\ctxsumTot\Gamma_2}{\PairExp{}{\ve_1}{\ve_2}{}}{\PairT{\tau_1}{\tau_2}}{\VarSet_1\cup\VarSet_2}
}{
}
\\[3ex]
\NamedRule{t-fun}{
\IsWFExpGr{\Gamma,\VarType{\x}{\tau_1}}{\e}{{\tau_2}}{\graph}
}{ 
\IsWFValGr{\Gamma}{\Fun{\x}{\e}}{\funTypeGr{\tau_1}{\tau_2}{\x}{\graph}} {\emptyset}
}
{
}
\\[3ex]
\end{array}
\end{math}
\hrule
\begin{math}
\begin{array}{l}
\\
\NamedRule{t-ret}{\IsWFValGr{\Gamma}{\ve}{\tau}{\VarSet}}{\IsWFExpGr{\Gamma}{\Return{\ve}}{\tau}{\sum_{\x\in\VarSet} \Edge{\x}{\res}}}{}\BigSpace
\\[3ex]
\NamedRule{t-let}{
\IsWFExpGr{\Gamma_1}{\e_1}{\tau_1}{\graph_1}\BigSpace\IsWFExpGr{\Gamma_2,\VarType{\x}{\tau_1}}{\e_2}{\tau_2}{\graph_2}
}{\IsWFExpGr{\Gamma_1\ctxsumTot\Gamma_2}{\Let{\x}{\e_1}{\e_2}}{\tau_2}{\RemoveGr{(\SumGr{\SubstGr{\graph_1}{\x}{\res}}{\graph_2})}{\{\x\}}}
}{
}
\\[3ex]
\NamedRule{t-match-p}{
  \IsWFValGr{\Gamma_1}{\ve}{\PairT{\tau_1}{\tau_2}}{\VarSet}
  \Space
  \IsWFExpGr{\Gamma_2,\VarType{\x}{\tau_1},\VarType{\y}{\tau_2}}{\e}{\tau}{\graph}
}{ \IsWFExpGr{\Gamma_1\ctxsumTot\Gamma_2}{\Match{\x}{\y}{\ve}{\e}}{\tau}{\RemoveGr{(\graph\cup\bigcup_{\z\in\VarSet}\{(\x,\z),(\y,\z)\})}{\{\x,\y\}}}}
{ 
}
\\[3ex]
\NamedRule{t-app}{
  \IsWFValGr{\Gamma_1}{\ve_1}{\funTypeGr{\tau_1}{\tau_2}{\x}{\graph}}{\_}
  \Space  
  \IsWFValGr{\Gamma_2}{\ve_2}{\tau_1}{\VarSet}
}{ \IsWFExpGr{\Gamma_1\ctxsumTot\Gamma_2}{\App{\ve_1}{\ve_2}}{\tau_2}{\SubstGr{\graph}{\VarSet}{\x}} }
{
}
\\[3ex]
\NamedRule{t-upd}{
\IsWFValGr{\Gamma_1}{\ve_1}{{\RefTV{\VarSet}{\tau}}}{\VarSet_1}\BigSpace
\IsWFValGr{\Gamma_2}{\ve_2}{{\tau}}{\VarSet_2}
}{
\IsWFExpGr{\Gamma_1\ctxsumTot\Gamma_2}{\Assign{\ve_1}{\ve_2}}{\tau}{\bigcup_{\x\in\VarSet_1,y\in\VarSet_2}\{(\x,\y)\}\cup\bigcup_{\z\in\VarSet}\{(\z,\res)\}}
}
{
\VarSet_2\subseteq\VarSet
}
\\[3ex]
\NamedRule{t-ref}{
\IsWFValGr{\Gamma}{\ve}{\tau}\VarSet
}
{
\IsWFExpGr{\Gamma}{\RefExp{\ve}}{\RefTV{\VarSet}{\tau}}{\emptygraph}
}
{}
\Space
\NamedRule{t-deref}{
\IsWFValGr{\Gamma}{\ve}{\RefTV{\VarSet}{\tau}}{\VarSet'}
}{
\IsWFExpGr{\Gamma}{\Deref{\ve}}{\tau}{\sum_{\x\in\VarSet\cup\VarSet'} \Edge{\x}{\res}}
}
{}
\\[3ex]
\NamedRule{t-free}{
\IsWFValGr{\Gamma}{\ve}{\RefTV{\VarSet}{\tau}}{\_}
}{
\IsWFExpGr{\Gamma}{\FreeExp{\ve}}{\tau}{\brb\sum_{\x\in\VarSet} \Edge{\x}{\res}\erb}
}
{}
\\[4ex]
\end{array}
\end{math}
\caption{Effect system}
\label{fig:typing} 

\end{figure}



